% Context from chapters-tex/denoising.tex; for mathematical reference, not compilation.
rvation~\cite{donoho-shrinkage}.

\myfigure{
\tabquatre{
\image{denoising}{.24}{wavthresh-noisy}&
\image{denoising}{.24}{wavthresh-coefs-noisy}&
\image{denoising}{.24}{wavthresh-coef-denoised}&
\image{denoising}{.24}{wavthresh-denoised}\\
$f$ & $\dotp{f}{\psi_{j,n}^\om}$ & $S_T(\dotp{f}{\psi_{j,n}^\om})$ & $\tilde f$
}
}{ 
	Denoising using thresholding of wavelet coefficients.   
}{fig-wavthresh}

Hard thresholding retains only coefficients whose magnitude exceeds the threshold $T$:
\eq{
		\tilde f = \sum_{ |\dotp{f}{\psi_m}|>T } \dotp{f}{\psi_m} \psi_m
		 = \sum_{ m } S_T(\dotp{f}{\psi_m}) \psi_m.
}
where $S_T$ is defined in~\eqref{eq-hard-thresh-approx}. Retaining the $M$ largest noisy coefficients gives the approximation $\tilde f = f_M$ of $f$. Figure \ref{fig-wavthresh} illustrates how a suitable $T$ removes much of the noise while preserving sharp features.

                                                                                                             

                                      
                                                                                                                                                                              

                                                     
\subsection{Soft Thresholding}

We recall that the hard thresholding operator is defined as
\eql{\label{eq-hard-thresh-denoise}
	S_T(x) = S_T^0(x) = \choice{ x \qifq |x|>T, \\ 0 \qifq |x|\leq T. }
}
Hard thresholding makes a discontinuous keep-or-discard decision, which can create artifacts. Soft thresholding is a continuous alternative:
\eql{\label{eq-soft-thresh-denoise}
	 S_T^1( x ) = \max( 1-T/|x|, 0 ) x.
}
Set $S_T^1(0)=0$ by continuity. Figure~\ref{fig-thresholding-hard-vs-soft} compares the two nonlinear maps.

\myfigure{
\BookDrawing{tikz/denoising-hard-soft-threshold}
}{ 
	Hard and soft thresholding functions.   
}{fig-thresholding-hard-vs-soft}

For $q=0$ and $q=1$, these thresholding rules define two different estimators
\eql{\label{eq-hard-soft-thresh}
	\tilde f^q = \sum_m S_T^q( \dotp{f}{\psi_m} ) \psi_m
}

\myfigure{
\image{denoising}{.6}{wavthresh-2d-hard-vs-soft}
}{ 
	SNR as a function of $T/\si$ for hard and soft thresholding. 
}{fig-wavthresh-2d-hard-vs-soft}


  
\paragraph{Preserving coarse-scale coefficients.}

Soft thresholding $S_T^1$ biases retained coefficients by reducing their magnitudes.
Shrinking coarse wavelet coefficients can introduce unwanted low-frequency artifacts. At a coarsest scale $2^{j_0}$, it is therefore common to retain the approximation coefficients unchanged. In one dimension, this gives
\eq{
	\tilde f^1 = \sum_{0 \leq n < 2^{-j_0}} \dotp{f}{\phi_{j_0,n}} \phi_{j_0,n}
	+ \sum_{j=J+1}^{j_0}\sum_{0\leq n<2^{-j}}S_T^1(\dotp{f}{\psi_{j,n}})\psi_{j,n}.
}
                                                                                                 

                                       
                                                                                                                                                                               


  
\paragraph{Empirical choice of the threshold.}

Figure \ref{fig-wavthresh-2d-hard-vs-soft} plots SNR against the threshold $T$ for a fixed natural image $f_0$. In this experiment, hard thresholding performs best near $T \approx 3\si$, and soft thresholding near $T\approx 3\si/2$. The best soft-thresholded estimate has a slightly higher SNR.

\myfigure{
\tabdeux{
\image{denoising}{.3}{wavthresh-2d-hard}&
\image{denoising}{.3}{wavthresh-2d-soft}\\
20.9dB & 21.8dB
}
}{ 
	Comparison of hard (left) and soft (right) thresholding.  
}{fig-wavthresh-2d}

These values describe the displayed experiments; the best threshold depends on the image, basis, noise level, and loss function.


                                                     
\subsection{Minimax Optimality of Thresholding}
\label{subsec-minimax-denoise}

  
\paragraph{Estimating sparse coefficients.}

To analyze the estimator and guide the choice of $T$, we first assume that the coefficients
\eq{
	a_{0,m}=\dotp{f_0}{\psi_m}\in\RR,\qquad a_0\in\RR^N
}
are sparse: most entries $a_{0,m}$ vanish, so the $\lzero$ count
\eq{
	\norm{a_0}_0 = \#\enscond{m}{a_{0,m} \neq 0}
}
is small. 
As shown in \eqref{eq-noisy-transformed}, noisy coefficients 
\eq{
	\dotp{f}{\psi_m} = a_m = a_{0,m} + z_m
}
are perturbed by additive Gaussian white noise of variance $\si^2$. Figure \ref{fig-sparse-signal} shows an example of such a noisy sparse signal.

\myfigure{
\image{denoising}{.45}{sparse-signal}
\image{denoising}{.45}{sparse-signal-noisy}
}{ 
	Left: clean sparse coefficients $a_0$; right: noisy coefficients $a$. 
}{fig-sparse-signal}


           
                                            
   
                                                     
                           

  
\paragraph{Universal threshold value.}

When all nonzero clean coefficients are well separated from the noise level, choosing a threshold comparable to the largest pure-noise coefficient suppresses false detections:
\eq{T\approx\tau_N,\qquad\tau_N=\max_{0\leq m<N}|z_m|.}
This motivates a deterministic threshold based on the typical size of $\tau_N$; it is not an exact minimization formula for the realized denoising error.
The random maximum $\tau_N$ depends on $N$. For fixed $\sigma$, its mean is asymptotic to $\sigma\sqrt{2\log N}$ as $N\to\infty$. Its variance tends to zero with $N$, so $\tau_N$ concentrates increasingly near its mean. Figure \ref{fig-max-gaussian-mean} illustrates these trends numerically.

\myfigure{
\image{denoising}{.6}{max-gaussian-mean}
\image{denoising}{.63}{max-gaussian-std-corrected}
}{ 
	Empirical mean (top) and standard deviation (bottom) of the largest absolute Gaussian noise coefficient. 
}{fig-max-gaussian-mean}


  
\paragraph{Asymptotic optimality.}

Donoho and Johnstone~\cite{donoho-shrinkage} showed that the universal threshold $T = \si \sqrt{2 \log(N)}$ gives a risk bound for signals that admit accurate nonlinear approximations in $\{\psi_m\}_m$. The risk is within a logarithmic factor of an ideal coefficient-selection risk.


\begin{thm}\label{thm-nonlinear-denoising}
Let $N\geq2$ and $Y=f_0+w$, with $w\sim\mathcal N(0,\sigma^2I_N)$, in a real orthonormal basis. Hard or soft thresholding~\eqref{eq-hard-soft-thresh} with
\eq{T=\sigma\sqrt{2\ln N}}
satisfies an oracle inequality of the form
\eq{
 \mathbb E\norm{f_0-\tilde f^q}^2
 \leq C_0(1+\ln N)\left[\sigma^2+\min_{0\leq M\leq N}\left(\norm{f_0-f_{0,M}^{\mathrm{nlin}}}^2+M\sigma^2\right)\right],
}
where $C_0$ is an absolute constant and $f_{0,M}^{\mathrm{nlin}}$ retains the $M$ largest clean coefficients. In particular, if $\norm{